\documentclass{article}
\author{Dawid S. Jakubczyk}
\usepackage{polski}
\usepackage[table]{xcolor}
\begin{document}
Dziedzina: Przychodnia medyczna\\
Obszar: Obsluga wizyt lacznie z zarzadzaniem pacjentami i lekarzami, zarzadzanie wyposazeniem przychodni, proces obslugi wizyt\\

Kontekst dziedziny problemowej:\\
Pacjent przychodzi do rejestracji rejestrujac sie na wizyte, przychodnia rezerwuje gabinet lekarski, czas lekarza, i jesli jest wymagany sprzet medyczny to rowniez jest przypisywany do wizyty. Po wizycie publicznej NFZ rozlicza sie z przychodnia, jesli wizyta byla prywatna pacjent reguluje naleznosci w rejestracji.
Przychodnia odpowiada za przeslanie informacji do systemu recept 


Procesy i aktorzy biznesowi:\\
\begin{tabular}{|m{6em}|m{10em}|m{15em}|m{10em}|}
\hline
\rowcolor{lightgray} Proces & Aktor Biznesowy & Funkcje/zadania & Dane\\
\hline
Rejestracja Pacjenta	& Pacjent,\newline pracownik rejestracji & Pracownik rejestracji ustala wizyte z lekarzem & Dane Pacjenta, Informacje medyczne\\
\hline
Wizyta lekarska & Pacjent,\newline Lekarz	& Pacjent konsultuje sie z Lekarzemz\newline Lekarz wystawia recepty albo skierowanie do specjalisty& Informacje medyczne\\
\hline
Wystawienie recept & Krajowa Baza Danych recept&Lekarz wystawia recepte i informacja jest przesylana do Bazy Danych&Informacje o Leku, nr Lekarza\\
\hline
\end{tabular}




.
\end{document}